% GENERATED FILE. Edit canonical lessons, then run npm run generate:books.
% canonical-source-hash: fnv1a64:f5666a7b5c5b899a
% canonical-lessons: FR-C461-bruine, FR-C461-bruiner, FR-C461-averse, FR-C461-eclaircie, FR-C461-degager

\chapter{Bruine et averses}
\label{ch:fr-bruine-et-averses}
\hlchaptermodality{461}

\begin{chapteropening}
\textbf{By the end of this chapter:} \emph{I can talk about showers: drizzle, drizzling, the shower, the bright spell and the sky clearing.}

Complete the last lesson of chapter 461: I can talk about showers: drizzle, drizzling, the shower, the bright spell and the sky clearing.
\end{chapteropening}

\section[la bruine]{la bruine — the drizzle}

\label{lesson:FR-C461-bruine}



\begin{quote}
Before the new one: say the French for sunny, then the French for the thermometer.
\end{quote}

\subsection*{You'll want to know: la bruine}
\textbf{la bruine} — \textquotedblleft{}the drizzle\textquotedblright{}.

\begin{quote}
Une bruine fine tombe depuis ce matin.
\end{quote}

\begin{quote}
\emph{A fine drizzle has been falling since this morning.}
\end{quote}

Very light rain.

\begin{cousinweb}
\emph{bruine} is usually traced to Latin \emph{pruīna} 'hoarfrost' (Lewis \& Short), but the path is not settled, and EtymDB does not list it.
\end{cousinweb}

\begin{grammarlens}[title={What you've built}]
One more word for telling, explaining and describing at B1.
\end{grammarlens}

\subsection*{Your turn}
\begin{itemize}
\raggedright
  \item \emph{Say it:} \emph{la bruine}
  \item \emph{Say it:} \emph{la bruine}, once more
  \item \emph{Say it:} a sentence of your own with \emph{la bruine}
  \item \emph{Recall:} say the French for sunny, then the French for the thermometer, then say \emph{la bruine} again
\end{itemize}

\begin{tcolorbox}[breakable,colback=teal!4,colframe=teal!35!black,title={Before you move on}]
What does la bruine mean? (\textquotedblleft{}The drizzle\textquotedblright{}.) Say it once more.
\end{tcolorbox}
\section[bruiner]{bruiner — to drizzle}

\label{lesson:FR-C461-bruiner}



\begin{quote}
Before the new one: say the French for the thermometer, then the French for the drizzle.
\end{quote}

\subsection*{You'll want to know: bruiner}
\textbf{bruiner} — \textquotedblleft{}to drizzle\textquotedblright{}.

\begin{quote}
Il bruine, mais nous sortons quand même.
\end{quote}

\begin{quote}
\emph{It is drizzling, but we are going out anyway.}
\end{quote}

Only with \emph{il}: \emph{il bruine}.

\begin{cousinweb}
\emph{bruiner} is built on \emph{bruine} 'drizzle' (EtymDB lists the verb).
\end{cousinweb}

\begin{grammarlens}[title={What you've built}]
One more word for telling, explaining and describing at B1.
\end{grammarlens}

\subsection*{Your turn}
\begin{itemize}
\raggedright
  \item \emph{Say it:} \emph{bruiner}
  \item \emph{Say it:} \emph{bruiner}, once more
  \item \emph{Say it:} a sentence of your own with \emph{bruiner}
  \item \emph{Recall:} say the French for the thermometer, then the French for the drizzle, then say \emph{bruiner} again
\end{itemize}

\begin{tcolorbox}[breakable,colback=teal!4,colframe=teal!35!black,title={Before you move on}]
What does bruiner mean? (\textquotedblleft{}To drizzle\textquotedblright{}.) Say it once more.
\end{tcolorbox}
\section[l'averse]{l'averse — the rain shower, the downpour}

\label{lesson:FR-C461-averse}



\begin{quote}
Before the new one: say the French for the drizzle, then the French for to drizzle.
\end{quote}

\subsection*{You'll want to know: l'averse}
\textbf{l'averse} — \textquotedblleft{}the rain shower, the downpour\textquotedblright{}.

\begin{quote}
Nous avons attendu la fin de l'averse sous un arbre.
\end{quote}

\begin{quote}
\emph{We waited for the shower to end under a tree.}
\end{quote}

A feminine noun. Short and heavy.

\begin{cousinweb}
\emph{averse} comes from \emph{à verse} 'pouring down', from \emph{verser} 'to pour', which goes back to Latin \emph{versāre}, from \emph{vertere} 'to turn' (Lewis \& Short). EtymDB does not list it.
\end{cousinweb}

\begin{grammarlens}[title={What you've built}]
One more word for telling, explaining and describing at B1.
\end{grammarlens}

\subsection*{Your turn}
\begin{itemize}
\raggedright
  \item \emph{Say it:} \emph{l'averse}
  \item \emph{Say it:} \emph{l'averse}, once more
  \item \emph{Say it:} a sentence of your own with \emph{l'averse}
  \item \emph{Recall:} say the French for the drizzle, then the French for to drizzle, then say \emph{l'averse} again
\end{itemize}

\begin{tcolorbox}[breakable,colback=teal!4,colframe=teal!35!black,title={Before you move on}]
What does l'averse mean? (\textquotedblleft{}The rain shower, the downpour\textquotedblright{}.) Say it once more.
\end{tcolorbox}
\section[l'éclaircie]{l'éclaircie — the bright spell, the sunny interval}

\label{lesson:FR-C461-eclaircie}



\begin{quote}
Before the new one: say the French for to drizzle, then the French for the rain shower.
\end{quote}

\subsection*{You'll want to know: l'éclaircie}
\textbf{l'éclaircie} — \textquotedblleft{}the bright spell, the sunny interval\textquotedblright{}.

\begin{quote}
Après la pluie, il y a eu une belle éclaircie.
\end{quote}

\begin{quote}
\emph{After the rain there was a lovely bright spell.}
\end{quote}

A feminine noun. \emph{Des averses et des éclaircies} — sunshine and showers.

\begin{cousinweb}
\emph{éclaircie} is built on \emph{éclaircir} 'to brighten, to clear up', from \emph{clair} 'clear', which is inherited from Latin \emph{clārus} (Lewis \& Short). EtymDB does not list it.
\end{cousinweb}

\begin{grammarlens}[title={What you've built}]
One more word for telling, explaining and describing at B1.
\end{grammarlens}

\subsection*{Your turn}
\begin{itemize}
\raggedright
  \item \emph{Say it:} \emph{l'éclaircie}
  \item \emph{Say it:} \emph{l'éclaircie}, once more
  \item \emph{Say it:} a sentence of your own with \emph{l'éclaircie}
  \item \emph{Recall:} say the French for to drizzle, then the French for the rain shower, then say \emph{l'éclaircie} again
\end{itemize}

\begin{tcolorbox}[breakable,colback=teal!4,colframe=teal!35!black,title={Before you move on}]
What does l'éclaircie mean? (\textquotedblleft{}The bright spell, the sunny interval\textquotedblright{}.) Say it once more.
\end{tcolorbox}
\section[dégager]{dégager — to clear (of the sky); to free}

\label{lesson:FR-C461-degager}



\begin{quote}
Before the new one: say the French for the rain shower, then the French for the bright spell.
\end{quote}

\subsection*{You'll want to know: dégager}
\textbf{dégager} — \textquotedblleft{}to clear (of the sky); to free\textquotedblright{}.

\begin{quote}
Vers midi, le ciel s'est dégagé.
\end{quote}

\begin{quote}
\emph{Around midday the sky cleared.}
\end{quote}

\emph{Le ciel se dégage} — the sky is clearing; \emph{dégager le passage} — to clear the way.

\begin{cousinweb}
\emph{dégager} is built from \emph{dé-} and \emph{gage} 'pledge' (EtymDB lists the verb), which comes from a Frankish word: first to take back what was pledged, then to set free.
\end{cousinweb}

\begin{grammarlens}[title={What you've built}]
One more word for telling, explaining and describing at B1.
\end{grammarlens}

\subsection*{Your turn}
\begin{itemize}
\raggedright
  \item \emph{Say it:} \emph{dégager}
  \item \emph{Say it:} \emph{dégager}, once more
  \item \emph{Say it:} a sentence of your own with \emph{dégager}
  \item \emph{Recall:} say the French for the rain shower, then the French for the bright spell, then say \emph{dégager} again
\end{itemize}

\begin{tcolorbox}[breakable,colback=teal!4,colframe=teal!35!black,title={Before you move on}]
What does dégager mean? (\textquotedblleft{}To clear (of the sky); to free\textquotedblright{}.) Say it once more.
\end{tcolorbox}
